\section{Технико-экономическое обоснование затрат на разработку программного комплекса}

\subsection{Характеристика программного комплекса}

Разрабатываемый программный комплекс представляет собой асимметричную систему гибридного сжатия и восстановления цифровых изображений, интегрирующую методы дискретного вейвлет-преобразования (DWT) и легковесных моделей глубокого машинного обучения (ESPCN, FSRCNN). Актуальность разработки обусловлена необходимостью минимизации объемов передаваемых данных в сетях с ограниченной пропускной способностью периферийных устройств.

Система нейросетевого сжатия на основе вейвлет-декомпозиции предназначена для компактного представления мультимедийных данных при сохранении семантически значимых признаков. Потенциальные области применения результатов разработки: передача данных по каналам с ограниченной пропускной способностью (IoT, мобильные сети), архивное хранение больших графических наборов данных, а также предобработка визуальных признаков в составе комплексных систем машинного зрения.

Ключевым технологическим и экономическим преимуществом предложенного решения является реализация алгоритма тайловой маршрутизации и механизма субпиксельной свертки. Это устраняет необходимость закупки дорогостоящего специализированного оборудования с дискретными графическими ускорителями, позволяя осуществлять инференс (декодирование) на классических многоядерных центральных процессорах.

В данном разделе определяется экономическая целесообразность разработки, которая обоснована расчетом сметы затрат на выполнение проектных и программных работ ($\text{С}_1$). Все расчеты приведены в белорусских рублях по состоянию на \targetYear~год.

\subsection{Смета затрат на разработку программного комплекса}

Расчет затрат на оплату труда персонала ($\text{З}_{\text{о}}$) представлен в таблице~6.1.

\begin{table}[ht]
\caption{Расчет основной зарплаты исполнителей проектных работ}
\begin{tabularx}{\textwidth}{|X|C{0.16\textwidth}|C{0.18\textwidth}|C{0.22\textwidth}|C{0.14\textwidth}|}
\hline
\centering\arraybackslash Исполнители & \centering\arraybackslash Количество исполнителей, чел. & \centering\arraybackslash Трудоемкость, мес. & \centering\arraybackslash Средне\-месячная зарплата, руб.* & \centering\arraybackslash Сумма, руб. \\
\hline
Инженер-программист & 1 & 1,0 & 2500,00 & 2500,00 \\
\hline
Инженер по машинному обучению & 1 & 1,0 & 3000,00 & 3000,00 \\
\hline
\multicolumn{4}{|l|}{Премия (40\,\%)} & 2200,00 \\
\hline
\multicolumn{4}{|l|}{Всего основная заработная плата ($\text{З}_{\text{о}}$)} & 7700,00 \\
\hline
\multicolumn{5}{|p{\dimexpr\textwidth-2\tabcolsep-2\arrayrulewidth\relax}|}{*Примечание~--- статистика уровня заработных плат ИТ-специалистов взята на основе открытых данных портала rabota.by~\cite{rabota_by_2026} по состоянию на \targetYear~год.} \\
\hline
\end{tabularx}
\end{table}

Итоговая сумма основной заработной платы ($\text{З}_{\text{о}}$) составляет 7700,00~рублей.

Дополнительная зарплата ($\text{З}_{\text{д}}$) определяется по формуле:
\begin{equation}
    \text{З}_{\text{д}} = \frac{\text{З}_{\text{о}} \cdot \text{Н}_{\text{д}}}{100\,\%},
\end{equation}
\noindent где $\text{Н}_{\text{д}}$~--- норматив дополнительной заработной платы основных работников, $\text{Н}_{\text{д}} = 20\,\%$.

Подставляя значения, получаем $\text{З}_{\text{д}} = 7700 \cdot 0,2 = 1540,00$~руб.

Отчисления на социальные нужды ($\text{Р}_{\text{соц}}$) определяются по формуле:
\begin{equation}
    \text{Р}_{\text{соц}} = \frac{(\text{З}_{\text{о}} + \text{З}_{\text{д}}) \cdot \text{Н}_{\text{соц}}}{100\,\%},
\end{equation}
\noindent где $\text{Н}_{\text{соц}}$~--- страховые взносы на обязательное социальное страхование наемных работников (34\,\%) и обязательное страхование от несчастных случаев на производстве (0,6\,\%), итого 34,6\,\%.

Подставляя значения, получаем $\text{Р}_{\text{соц}} = (7700 + 1540) \cdot 0,346 = 3197,04$~руб.

Прочие расходы ($\text{Р}_{\text{пр}}$) определяются по формуле:
\begin{equation}
    \text{Р}_{\text{пр}} = \frac{\text{З}_{\text{о}} \cdot \text{Н}_{\text{пр}}}{100\,\%},
\end{equation}
\noindent где $\text{Н}_{\text{пр}}$~--- норматив прочих расходов (1,5\,\%).

Подставляя значения, получаем $\text{Р}_{\text{пр}} = 7700 \cdot 0,015 = 115,50$~руб.

Накладные расходы ($\text{Р}_{\text{накл}}$) определяются по формуле:
\begin{equation}
    \text{Р}_{\text{накл}} = \frac{\text{З}_{\text{о}} \cdot \text{Н}_{\text{накл}}}{100\,\%},
\end{equation}
\noindent где $\text{Н}_{\text{накл}}$~--- норматив накладных расходов (50\,\%).

Подставляя значения, получаем $\text{Р}_{\text{накл}} = 7700 \cdot 0,5 = 3850,00$~руб.

Полная себестоимость ($\text{С}_{\text{п}}$) рассчитывается как сумма всех статей затрат:
\begin{equation}
    \text{С}_{\text{п}} = \text{З}_{\text{о}} + \text{З}_{\text{д}} + \text{Р}_{\text{соц}} + \text{Р}_{\text{пр}} + \text{Р}_{\text{накл}}.
\end{equation}

Подставляя значения, получаем $\text{С}_{\text{п}} = 7700 + 1540 + 3197,04 + 115,50 + 3850 = 16~402,54$~руб.

Плановая прибыль ($\text{П}_{\text{п}}$) определяется следующим образом:
\begin{equation}
    \text{П}_{\text{п}} = \frac{\text{С}_{\text{п}} \cdot \text{Р}_{\text{п}}}{100\,\%},
\end{equation}
\noindent где $\text{Р}_{\text{п}}$~--- уровень рентабельности (20\,\%).

Подставляя значения, получаем $\text{П}_{\text{п}} = 16~402,54 \cdot 0,2 = 3280,51$~руб.

Цена предприятия ($\text{Ц}_{\text{пр}}$) определяется по формуле:
\begin{equation}
    \text{Ц}_{\text{пр}} = \text{С}_{\text{п}} + \text{П}_{\text{п}}.
\end{equation}

Подставляя значения, получаем $\text{Ц}_{\text{пр}} = 16~402,54 + 3280,51 = 19~683,05$~руб.

Налог на добавленную стоимость (НДС) определяется следующим образом:
\begin{equation}
    \text{НДС} = \frac{\text{Ц}_{\text{пр}} \cdot \text{Н}_{\text{дс}}}{100\,\%},
\end{equation}
\noindent где $\text{Н}_{\text{дс}}$~--- ставка налога на добавленную стоимость (20\,\%).

Подставляя значения, получаем $\text{НДС} = 19~683,05 \cdot 0,2 = 3936,61$~руб.

Итоговая смета на разработку ($\text{С}_1$):
\begin{equation}
    \text{С}_1 = \text{С}_{\text{п}} + \text{П}_{\text{п}} + \text{НДС}.
\end{equation}

Подставляя значения, получаем $\text{С}_1 = 16~402,54 + 3280,51 + 3936,61 = 23~619,66$~руб.

Итоговые результаты расчетов сведены в таблицу~6.2.

\begin{table}[ht]
\caption{Расчет сметы затрат на разработку}
\begin{tabularx}{\textwidth}{|X|c|C{0.25\textwidth}|}
\hline
\centering\arraybackslash Наименование статьи затрат & \centering\arraybackslash Условные обозначения & \centering\arraybackslash Значение, руб. \\
\hline
1 Заработная плата & & \\
1.1 Основная заработная плата & $\text{З}_{\text{о}}$ & 7700,00 \\
1.2 Дополнительная заработная плата & $\text{З}_{\text{д}}$ & 1540,00 \\
\hline
2 Отчисления в фонд социальной защиты & $\text{Р}_{\text{соц}}$ & 3197,04 \\
\hline
3 Прочие расходы & $\text{Р}_{\text{пр}}$ & 115,50 \\
\hline
4 Накладные расходы & $\text{Р}_{\text{накл}}$ & 3850,00 \\
\hline
5 Полная себестоимость & $\text{С}_{\text{п}}$ & 16~402,54 \\
\hline
6 Плановая прибыль & $\text{П}_{\text{п}}$ & 3280,51 \\
\hline
7 Цена предприятия & $\text{Ц}_{\text{пр}}$ & 19~683,05 \\
\hline
8 Налог на добавленную стоимость & $\text{НДС}$ & 3936,61 \\
\hline
9 Итоговая смета на разработку & $\text{С}_1$ & 23~619,66 \\
\hline
\end{tabularx}
\end{table}

Таким образом, экономический эффект от реализации гибридного алгоритма сжатия изображений равен плановой прибыли и составляет 3280,51~руб. Оптимизация архитектуры под центральный процессор дополнительно позволяет сэкономить на капитальных затратах при внедрении (отсутствует необходимость закупки дорогостоящих графических видеокарт).